% Motivationsschreiben DI Alexander Sommer - Bewerbung USTP St. Pölten
% Position: Departmentleiter*in / Institutsleiter*in Digital Business & Innovation (REC-0062)
% Compiler: Tectonic / pdfLaTeX
\documentclass[10.5pt]{article}
\usepackage[a4paper,top=16mm,bottom=15mm,left=21mm,right=21mm]{geometry}
\usepackage[T1]{fontenc}
\usepackage[utf8]{inputenc}
\usepackage[ngerman]{babel}
\IfFileExists{sourcesanspro.sty}%
  {\usepackage[default,semibold]{sourcesanspro}}%
  {\usepackage[scaled=0.92]{helvet}\renewcommand{\familydefault}{\sfdefault}}
\usepackage{microtype}
\usepackage{xcolor}
\usepackage[hidelinks]{hyperref}

\definecolor{accent}{HTML}{1F3A5F}
\definecolor{brandblue}{HTML}{1E3A8A}
\definecolor{muted}{HTML}{5A6472}

\hypersetup{
  pdftitle={Motivationsschreiben DI Alexander Sommer - USTP St. Pölten},
  pdfauthor={DI Alexander Sommer}
}

\setlength{\parindent}{0pt}
\setlength{\parskip}{4.5pt}
\linespread{1.04}
\pagestyle{empty}

\begin{document}

% ---------------------------------------------------------------- Absender / Empfänger
\begin{minipage}[t]{0.56\textwidth}
  \textbf{\color{accent}\large DI Alexander Sommer}\\[1pt]
  Gregorygasse 12/6/2 · 1230 Wien\\[1pt]
  Tel.: +43 681 102 104 17 · \href{mailto:alexander.sommer@gmail.com}{alexander.sommer@gmail.com}\\[1pt]
  \href{https://alexander-sommer.vercel.app}{alexander-sommer.vercel.app}
\end{minipage}\hfill
\begin{minipage}[t]{0.42\textwidth}\raggedleft
  \textbf{USTP – University of Applied Sciences St. Pölten}\\
  Personalmanagement\\
  z.\,H. Frau Sonja Gutschi\\
  Campus-Platz 1 · 3100 St. Pölten
\end{minipage}

\vspace{8pt}
\hfill Wien, 23. Oktober 2026

\vspace{2pt}
{\large\bfseries\color{brandblue}Bewerbung als Departmentleiter / Institutsleiter Digital Business \& Innovation (w/m/d)}
\vspace{1pt}
\newline{\footnotesize\color{muted}\textbf{Referenzcode: REC-0062} · 35--40 Wochenstunden · Lehre \& Forschung}

\vspace{4pt}
Sehr geehrte Frau Gutschi, sehr geehrte Mitglieder der Auswahlkommission,

die University of Applied Sciences St. Pölten hat sich durch ihr dynamisches Wachstum, ihre Forschungsstärke und die Vorreiterrolle in der europäischen Hochschulallianz \textbf{E³UDRES²} als einer der profiliertesten Innovationsstandorte im tertiären Bildungsraum etabliert. Die Möglichkeit, das Department \emph{Digital Business \& Innovation} sowie das Forschungsinstitut strategisch wie operativ zu führen, zukunftsweisende Schwerpunkte rund um Künstliche Intelligenz zu setzen und Forschung, Lehre und Drittmittelakquise eng mit regionalen und internationalen Partnern zu verzahnen, begeistert mich außerordentlich.

Als Diplom-Ingenieur der Wirtschaftsinformatik (TU Wien) verbinde ich eine fundierte wissenschaftliche Qualifikation mit mehr als 15 Jahren nachweisbarer Führungs- und Budgetverantwortung im regulierten Großumfeld:

\begin{itemize}\setlength{\itemsep}{1.5pt}\setlength{\parskip}{0pt}
  \item \textbf{Führung \& Department-Steuerung:} Bei der Austrian Power Grid AG (APG), Wiener Wohnen und der OMV verantwortete ich geschäftskritische IT-Portfolios mit bis zu 85 Kernsystemen und Budgets in Millionenhöhe. Führungsarbeit verstehe ich als werteorientiertes Ermöglichen: In einer externen Management-Potenzialeinschätzung (PEP / IMDE) wurde mein Führungsprofil im oberen Dezil (Top 10\,\%) vergleichbarer Führungskräfte verortet. Ich führe wertschätzend, zielorientiert und fördere Diversität sowie Chancengleichheit aus tiefer Überzeugung.

  \item \textbf{Wissenschaftliche Exzellenz \& Drittmittelerfahrung:} Mein Studium an der TU Wien schloss ich mit Bestnote ab (Publikation in den \emph{Forschungsmemoranden} des Instituts für Finanzwissenschaft). Ich publizierte peer-reviewed im \emph{International Journal of Pervasive Computing and Communications} (Emerald Insight), referierte international (Paris) und gewann den 1.~Preis beim Ideenwettbewerb der TU Wien (\emph{„Technik für Menschen“}). Im COMET-K2-Zentrum \emph{evolaris next level} wirkte ich an geförderten F\&E-Projekten mit (darunter die im \emph{AMS Report 117} dokumentierte Plattform \emph{Sankt Onlein} für IKT-Inklusion). Als Miterfinder eines vom Österreichischen Patentamt erteilten Hardware-Patents (AT 519762, FFG-gefördert) kenne ich den gesamten F\&E-Zyklus von der Drittmittelakquise bis zum Technologie-Transfer.

  \item \textbf{Schwerpunkt KI, Zukunftstechnologien \& Digital Management:} Um Zukunftsthemen wie Generative und Agentic AI fundiert in Curricula und Forschungsprojekte einzubringen, absolvierte ich 2025 das Programm \emph{Applied Agentic AI} an der M.I.T. Professional Education und entwickle aktiv eigene KI-Forschungsprototypen. Kombiniert mit 7 Jahren kontinuierlicher Expertise als Gartner Seat Holder (Executive \& Technical) bringe ich tiefes Verständnis für globale Technologieradare, Vendor-Governance und digitale Geschäftsmodelle mit.

  \item \textbf{Lehre, Internationalisierung \& Digitaler Humanismus:} Lehre an einer Fachhochschule muss inspirieren, methodisch schärfen und anwendungsorientiert begeistern. Mein Leitgedanke ist der \emph{Digitale Humanismus}: Technologie daran zu messen, wie wirksam sie Menschen befähigt und nachhaltigen Mehrwert stiftet. Diesen Anspruch möchte ich gemeinsam mit dem Team leben und über E³UDRES² strategische europäische Forschungsanträge und Hochschulkooperationen aktiv ausbauen.
\end{itemize}

Ich freue mich darauf, meine Erfahrung, mein Netzwerk und meine Leidenschaft für zukunftsfähige Hochschulentwicklung in die USTP einzubringen, und stehe Ihnen für ein persönliches Gespräch sehr gerne zur Verfügung.

\vspace{6pt}
Mit freundlichen Grüßen\\[16pt]
\textbf{DI Alexander Sommer}

\end{document}
